% =====================================================================
%  Ingenieria de Requisitos (ISR-401) -- UTEQ
%  Facultad de Ciencias de la Ingenieria -- Carrera de Ingenieria de Software
%  PRUEBA PRACTICA -- UNIDAD IV: Validacion, Gestion, Herramientas y Estandares
%  PARALELO B -- Caso: Sistema de Reserva de Citas Medicas.
%  Modalidad: individual, en clase, con repositorio.
%  Documento autonomo compilable. Formato tipo libro, A4.
%  Compilacion:  pdflatex main -> bibtex main -> pdflatex main -> pdflatex main
% =====================================================================
\documentclass[11pt,a4paper,oneside]{article}

\usepackage[utf8]{inputenc}
\usepackage[T1]{fontenc}
% Nombres en espanol sin depender de babel-spanish (no instalado en todos los sistemas)
\renewcommand{\tablename}{Tabla}
\renewcommand{\figurename}{Figura}
\usepackage[scaled=0.92]{helvet}
\usepackage{textcomp}

\usepackage[a4paper,top=2.4cm,bottom=2.3cm,left=2.5cm,right=2.3cm,headheight=15pt]{geometry}

\usepackage{amsmath,amssymb}
\usepackage{graphicx}
\usepackage[table,dvipsnames]{xcolor}
\usepackage{array}
\usepackage{tabularx}
\usepackage{multirow}
\usepackage{colortbl}
\usepackage{booktabs}
\usepackage{enumitem}
\usepackage[protrusion=true,expansion=false]{microtype}
\usepackage{parskip}
\usepackage{titlesec}
\usepackage{fancyhdr}
\usepackage{caption}
\usepackage{pdflscape}
\usepackage[numbers,sort&compress]{natbib}
\usepackage[colorlinks=true,linkcolor=darkblue,citecolor=darkblue,urlcolor=midblue,
	pdftitle={Prueba Practica Unidad IV - Paralelo B - Ingenieria de Requisitos ISR-401},
	pdfauthor={Ing. Gleiston Guerrero - UTEQ}]{hyperref}
\usepackage[most]{tcolorbox}

% ── Paleta institucional ─────────────────────────────────────────────
\definecolor{darkblue}{HTML}{1A3A5C}
\definecolor{midblue}{HTML}{2E6DA4}
\definecolor{gold}{HTML}{C8A951}
\definecolor{lightgray}{HTML}{F2F2F2}
\definecolor{rowgray}{HTML}{EEEEEE}
\definecolor{softblue}{HTML}{EAF1F8}
\definecolor{softred}{HTML}{F7E7E7}
\definecolor{darkred}{HTML}{9C2A2A}
\definecolor{softgreen}{HTML}{E7F1E7}

% ── Interlineado y columnas ──────────────────────────────────────────
\linespread{1.05}
\setlength{\parskip}{5pt}
\setlength{\parindent}{0pt}
\renewcommand{\arraystretch}{1.25}
\newcolumntype{L}[1]{>{\raggedright\arraybackslash}p{#1}}
\newcolumntype{C}[1]{>{\centering\arraybackslash}p{#1}}
\newcolumntype{R}[1]{>{\raggedleft\arraybackslash}p{#1}}

% ── Encabezados ──────────────────────────────────────────────────────
\titleformat{\section}{\sffamily\large\bfseries\color{darkblue}}{\thesection}{0.6em}{}
	[\vspace{-5pt}{\color{gold!85}\rule{\linewidth}{1pt}}]
\titlespacing*{\section}{0pt}{14pt}{6pt}
\titleformat{\subsection}{\sffamily\normalsize\bfseries\color{midblue}}{\thesubsection}{0.5em}{}
\titlespacing*{\subsection}{0pt}{10pt}{3pt}

\pagestyle{fancy}
\fancyhf{}
\renewcommand{\headrulewidth}{0.4pt}
\fancyhead[L]{\footnotesize\sffamily\color{darkblue}Prueba Pr\'actica\, \textbullet\, Unidad IV\, \textbullet\, Paralelo B}
\fancyhead[R]{\footnotesize\sffamily\color{darkblue}Ingenier\'ia de Requisitos\, \textbullet\, ISR-401}
\fancyfoot[C]{\footnotesize\sffamily\color{darkblue}\thepage}

\captionsetup{font=small,labelfont={bf,sf,color=darkblue},labelsep=period,skip=5pt}

% ── Cajas ────────────────────────────────────────────────────────────
\newtcolorbox{concepto}[1]{colback=softblue,colframe=darkblue,boxrule=0.6pt,arc=2.5pt,
	left=7pt,right=7pt,top=5pt,bottom=5pt,fonttitle=\bfseries\sffamily\color{white},
	coltitle=white,colbacktitle=darkblue,title={#1},breakable}
\newtcolorbox{piso}[1]{colback=softred,colframe=darkred,boxrule=1pt,arc=2.5pt,
	left=7pt,right=7pt,top=5pt,bottom=5pt,fonttitle=\bfseries\sffamily\color{white},
	coltitle=white,colbacktitle=darkred,title={#1},breakable}
\newtcolorbox{tarea}[1]{colback=white,colframe=midblue,boxrule=0.7pt,arc=2pt,
	left=7pt,right=7pt,top=5pt,bottom=5pt,fonttitle=\bfseries\sffamily\color{white},
	coltitle=white,colbacktitle=midblue,title={#1},breakable}

\newcommand{\hd}[1]{\cellcolor{darkblue}\textbf{\textcolor{white}{\small\sffamily #1}}}
\newcommand{\hm}[1]{\cellcolor{midblue}\textbf{\textcolor{white}{\small\sffamily #1}}}
\newcommand{\hg}[1]{\cellcolor{lightgray}\textbf{\color{darkblue}\small\sffamily #1}}

% =====================================================================
\begin{document}
\sffamily

% ── Portada breve / cabecera del instrumento ─────────────────────────
\begin{center}
	{\Large\bfseries\color{darkblue} UNIVERSIDAD T\'ECNICA ESTATAL DE QUEVEDO}\\[2pt]
	{\normalsize\color{midblue} Facultad de Ciencias de la Ingenier\'ia \textbullet\ Carrera de Ingenier\'ia de Software}\\[6pt]
	{\color{gold}\rule{0.9\linewidth}{1.2pt}}\\[6pt]
	{\large\bfseries\color{darkblue} PRUEBA PR\'ACTICA --- UNIDAD IV \textbullet\ PARALELO B}\\[2pt]
	{\normalsize\color{darkblue} Validaci\'on, Gesti\'on, Herramientas y Est\'andares de Requisitos}\\[2pt]
	{\small\color{black} Asignatura: Ingenier\'ia de Requisitos (ISR-401) \quad\textbullet\quad Docente: Ing. Gleiston Guerrero, Mg.}
\end{center}

\vspace{2pt}
\begin{center}\small
\begin{tabularx}{\linewidth}{@{}L{2.6cm} X L{2.6cm} X@{}}
	\hg{Estudiante} & ZAMBRANO MOYA ANGELO PAUL & \hg{Fecha} & 12/08/2026 \\
	\hg{Modalidad} & Individual, en clase & \hg{Duraci\'on} & 120 minutos \\
	\hg{Caso} & Sistema de Reserva de Citas M\'edicas & \hg{Ponderaci\'on} & 100 puntos (30 + 70) \\
\end{tabularx}
\end{center}

\vspace{3pt}
\begin{center}
\small\textbf{Repositorio:} \url{https://github.com/AngeloP2114/Evaluacion_Unidad_IV_Practica}
\end{center}

\vspace{2pt}
\begin{center}
\begin{minipage}[t]{0.98\linewidth}
\centering
\textbf{Captura del resumen del cuestionario}\par\vspace{2pt}
\includegraphics[width=\linewidth]{evidencias/resumen_cuestionario_SGA.png}
\end{minipage}

\vspace{4pt}
\begin{minipage}[t]{0.48\linewidth}
\centering
\textbf{Captura de la evaluaci\'on del cuestionario}\par\vspace{2pt}
\includegraphics[width=\linewidth,height=9.0cm,keepaspectratio]{evidencias/revision_intento_SGA.png}
\end{minipage}
\end{center}

\newpage

% ── Resultado de aprendizaje ─────────────────────────────────────────
\begin{concepto}{Resultado de aprendizaje evaluado}
El estudiante \textbf{construye y valida} los modelos de an\'alisis (datos, funci\'on y comportamiento) de un sistema, \textbf{especifica y gestiona} sus requisitos aplicando est\'andares (\textit{ISO/IEC/IEEE} 29148:2018 e \textit{ISO/IEC} 25010:2023) y \textbf{demuestra trazabilidad, control de cambios y gesti\'on de configuraci\'on} sobre un repositorio reproducible. Esta prueba eval\'ua \emph{desempe\~no pr\'actico}: no contiene preguntas te\'oricas; se califica lo que el estudiante \emph{produce}.
\end{concepto}

% #####################################################################
\section{Instrucciones generales y entregables}

La prueba se desarrolla de forma individual durante la sesi\'on. Todo el trabajo se versiona en un \textbf{repositorio Git p\'ublico} (GitHub o GitLab) creado por el estudiante. El entregable se redacta en \textbf{LaTeX} y se compila a \textbf{PDF} dentro del propio repositorio. Al finalizar, el estudiante sube al LMS (SGA) un \textbf{PDF con car\'atula} que contiene, \textbf{en una sola l\'inea}, la \textbf{URL} del repositorio.

\begin{enumerate}[leftmargin=1.4em,itemsep=1pt,topsep=2pt]
	\item \textbf{Repositorio} con: archivo principal \texttt{.tex}, archivo \texttt{referencias.bib} (si aplica), figuras, \texttt{PDF} compilado y \texttt{README.md}.
	\item \textbf{README.md} con las instrucciones exactas de compilaci\'on: compilador (\texttt{pdflatex}), orden de comandos, archivo principal y dependencias.
	\item \textbf{Car\'atula del PDF} con datos de identificaci\'on y la URL del repositorio en una sola l\'inea, m\'as la \textbf{captura del resumen} y la \textbf{captura de la evaluaci\'on} del cuestionario rendido en el SGA.
	\item Los modelos (clases, actividades, m\'aquina de estados) pueden dibujarse en \textbf{TikZ}, en una herramienta CASE (p.\,ej. Enterprise Architect / Sparx EA) o a mano y adjuntarse como imagen legible.
\end{enumerate}

% ── Criterios de piso ────────────────────────────────────────────────
\begin{piso}{Criterios de piso --- su incumplimiento implica calificaci\'on CERO (0) en toda la prueba}
\textbf{1. Evidencia del repositorio.} El estudiante debe subir al LMS (SGA) un documento \textbf{PDF con car\'atula} (datos de identificaci\'on) que muestre claramente, \textbf{en una sola l\'inea}, la \textbf{URL del repositorio} donde se encuentra el trabajo desarrollado. Si la URL est\'a rota, el repositorio no es accesible, o no se cumple esta directriz, la calificaci\'on es \textbf{CERO}.\\[3pt]
\textbf{2. Reproducibilidad del PDF.} Si el PDF \textbf{no se genera} correctamente a partir del LaTeX mediante \textbf{clonaci\'on del repositorio y compilaci\'on} del archivo \texttt{.tex}, la calificaci\'on es \textbf{CERO}. Las instrucciones de compilaci\'on (compilador, orden de comandos, archivo principal, dependencias) deben estar documentadas en el archivo \texttt{README.md} del repositorio.
\end{piso}

% #####################################################################
\section{Caso pr\'actico: Sistema de Reserva de Citas M\'edicas}

Una cl\'inica requiere un m\'odulo para gestionar las reservas de citas de sus pacientes. Un \textbf{Paciente} (identificado por c\'edula, nombre y correo) puede solicitar cero o m\'as \textbf{Citas}. Cada \textbf{Cita} (identificador, fecha-hora y estado) corresponde a un \textbf{M\'edico} (c\'odigo, nombre, especialidad) sobre un \textbf{Cupo} de su agenda (fecha-hora, disponibilidad). Al agendar una cita, el sistema verifica la disponibilidad de \textbf{cupo} en la agenda del m\'edico: si hay cupo, \emph{confirma} la cita; si no, \emph{notifica la no disponibilidad} y propone otro horario. Una cita confirmada puede \emph{atenderse}, o bien \emph{cancelarse} mientras no haya sido atendida; si el paciente no se presenta, la cita se marca como \emph{no asistida}. El sistema debe registrar la cita con rapidez y proteger los datos personales del paciente.

Sobre este caso, el estudiante desarrolla las diez actividades pr\'acticas siguientes. Todas las actividades usan el \emph{mismo} dominio, de modo que los modelos deben ser mutuamente consistentes.

% #####################################################################
\section{Actividades pr\'acticas (desarrollo --- 70 puntos)}

\begin{tarea}{P1. Modelo de datos --- Diagrama de clases UML \hfill (8 pts)}
Construya el \textbf{diagrama de clases UML} del dominio con al menos las clases \textsf{Paciente}, \textsf{Cita}, \textsf{M\'edico} y \textsf{Cupo} (agenda). Cada clase debe incluir sus \textbf{atributos} con tipo, al menos una \textbf{operaci\'on} pertinente y las \textbf{multiplicidades} en los extremos de cada asociaci\'on. Subraye o marque los identificadores.
\end{tarea}

\begin{center}
\includegraphics[width=0.98\linewidth]{figuras/P1_diagrama_clases.png}
\end{center}

\newpage


\begin{tarea}{P2. Modelo funcional --- Diagrama de actividades UML \hfill (8 pts)}
Modele el proceso \textbf{``Agendar cita''} con un \textbf{diagrama de actividades UML} que incluya: nodo inicial, al menos una \textbf{decisi\'on} (\textsf{\textquestiondown cupo disponible?}), los dos caminos alternativos (\emph{confirmar} / \emph{notificar no disponibilidad}), su convergencia y el nodo final. Use carriles de responsabilidad si distingue actores.
\end{tarea}

\begin{center}
\includegraphics[width=0.98\linewidth]{figuras/P2_diagrama_actividades.png}
\end{center}

\newpage


\begin{tarea}{P3. Modelo de comportamiento --- M\'aquina de estados UML \hfill (8 pts)}
Elabore la \textbf{m\'aquina de estados UML} del ciclo de vida de una \textsf{Cita} (p.\,ej. \textsf{Solicitada}, \textsf{Confirmada}, \textsf{Atendida}, \textsf{Cancelada}, \textsf{No asistida}). Etiquete cada \textbf{transici\'on} con su \textbf{evento}, incluya al menos una \textbf{condici\'on de guarda} y marque el estado final de aceptaci\'on. Se valora el uso de un superestado (statechart) para factorizar la transici\'on \emph{cancelar}.
\end{tarea}

\begin{center}
\includegraphics[width=0.98\linewidth]{figuras/P3_maquina_estados.png}
\end{center}

\newpage


\begin{tarea}{P4. Consistencia entre las tres perspectivas \hfill (7 pts)}
Elabore una \textbf{tabla de verificaci\'on cruzada} que relacione cada \textbf{evento} de P3 con la \textbf{funci\'on/acci\'on} de P2 que lo realiza y con la \textbf{entidad o atributo} de P1 que modifica. \textbf{Detecte al menos una inconsistencia} entre los modelos (evento sin funci\'on, entidad sin clase, etc.), \textbf{documente} el defecto y \textbf{corrija} el modelo afectado.
\end{tarea}

\subsection*{Tabla de verificaci\'on cruzada}
\small
\begin{tabularx}{\linewidth}{@{}L{3.0cm} X X@{}}
\toprule
\hd{Evento de P3} & \hd{Acci\'on o funci\'on de P2} & \hd{Entidad / atributo de P1 que modifica}\\
\midrule
\texttt{solicitarCita} & Solicitar cita / Registrar solicitud de cita & \texttt{Cita.idCita}, \texttt{Cita.fechaHora}, \texttt{Cita.estado}\\
\rowcolor{rowgray}
\texttt{verificarCupo} & Verificar disponibilidad del cupo / Consultar agenda del m\'edico & \texttt{Cupo.disponibilidad}\\
\texttt{confirmarCita} & Confirmar cita & \texttt{Cita.estado}, asociaci\'on Cita--Cupo\\
\rowcolor{rowgray}
\texttt{notificarNoDisponibilidad} & Notificar no disponibilidad & No modifica atributos; consulta \texttt{Cupo.disponibilidad}\\
\texttt{proponerHorario} & Proponer otro horario & \texttt{Cupo.fechaHora}\\
\rowcolor{rowgray}
\texttt{cancelarCita} & \textbf{No representado en P2} & \texttt{Cita.estado}\\
\texttt{atenderCita} & \textbf{No representado en P2} & \texttt{Cita.estado}\\
\rowcolor{rowgray}
\texttt{marcarNoAsistida} & \textbf{No representado en P2} & \texttt{Cita.estado}\\
\bottomrule
\end{tabularx}
\normalsize

\subsection*{Inconsistencia detectada}
\textbf{INC-01:} En la m\'aquina de estados P3 aparecen los eventos \texttt{cancelarCita}, \texttt{atenderCita} y \texttt{marcarNoAsistida}, pero en el diagrama de actividades P2 no est\'an representadas las acciones correspondientes. Por lo tanto, existe una inconsistencia entre el modelo funcional y el modelo de comportamiento.

\subsection*{Correcci\'on del modelo afectado}
\textbf{Modelo corregido:} P2 --- Diagrama de actividades.\\
\textbf{Correcci\'on realizada:} Se deben agregar al diagrama de actividades las siguientes acciones y decisiones para que coincida con P3:
\begin{itemize}[leftmargin=1.6em,itemsep=1pt,topsep=2pt]
\item Esperar fecha de la cita.
\item \textquestiondown Cita cancelada?
\item Cancelar cita.
\item \textquestiondown Paciente se presenta?
\item Atender cita.
\item Marcar cita no asistida.
\end{itemize}

\begin{center}
\includegraphics[width=0.96\linewidth]{figuras/P4_P2_corregido.png}
\end{center}

Al realizar la verificaci\'on cruzada entre P1, P2 y P3, se detect\'o la inconsistencia \textbf{INC-01}, debido a que en la m\'aquina de estados se definieron los eventos \texttt{cancelarCita}, \texttt{atenderCita} y \texttt{marcarNoAsistida}, pero dichas acciones no estaban representadas en el diagrama de actividades. Para corregir esta inconsistencia, se actualiz\'o el modelo funcional incorporando las acciones \emph{Cancelar cita}, \emph{Atender cita} y \emph{Marcar cita no asistida}, as\'i como las decisiones \emph{\textquestiondown Cita cancelada?} y \emph{\textquestiondown Paciente se presenta?}. Luego de la correcci\'on, los tres modelos mantienen coherencia entre eventos, funciones y entidades.

\newpage


\begin{tarea}{P5. Especificaci\'on de requisitos con esquema de atributos \hfill (10 pts)}
Redacte \textbf{4 requisitos}: \textbf{2 funcionales} y \textbf{2 no funcionales}. Cada requisito no funcional debe \textbf{nombrar una caracter\'istica de \textit{ISO/IEC} 25010:2023} (p.\,ej. Eficiencia de desempe\~no, Seguridad), fijar una \textbf{m\'etrica} y un \textbf{umbral} (p.\,ej. ``agendar la cita en $<$ 3 s''). Para cada requisito complete un \textbf{esquema de atributos}: identificador, prioridad, estado, fuente, estabilidad, versi\'on y m\'etodo de verificaci\'on.
\end{tarea}

\subsection*{RF-01 --- Registrar cita}
\begin{tabularx}{\linewidth}{@{}L{4.0cm} X@{}}
\toprule
\hd{Atributo} & \hd{Valor}\\
\midrule
Identificador & RF-01\\
\rowcolor{rowgray} Tipo & Funcional\\
Nombre & Registrar cita\\
\rowcolor{rowgray} Requisito & El sistema deber\'a registrar una cita solicitada por un paciente, asoci\'andola con un m\'edico y un cupo de su agenda.\\
Prioridad & Alta\\
\rowcolor{rowgray} Estado & En revisi\'on\\
Fuente & Caso pr\'actico: Sistema de Reserva de Citas M\'edicas\\
\rowcolor{rowgray} Estabilidad & Alta\\
Versi\'on & 1.0\\
\rowcolor{rowgray} M\'etodo de verificaci\'on & Prueba de aceptaci\'on PA-01\\
\bottomrule
\end{tabularx}

\subsection*{RF-02 --- Verificar disponibilidad del cupo}
\begin{tabularx}{\linewidth}{@{}L{4.0cm} X@{}}
\toprule
\hd{Atributo} & \hd{Valor}\\
\midrule
Identificador & RF-02\\
\rowcolor{rowgray} Tipo & Funcional\\
Nombre & Verificar disponibilidad del cupo\\
\rowcolor{rowgray} Requisito & El sistema deber\'a verificar la disponibilidad del cupo seleccionado antes de confirmar una cita.\\
Prioridad & Alta\\
\rowcolor{rowgray} Estado & En revisi\'on\\
Fuente & Caso pr\'actico: disponibilidad de cupos\\
\rowcolor{rowgray} Estabilidad & Alta\\
Versi\'on & 1.0\\
\rowcolor{rowgray} M\'etodo de verificaci\'on & Prueba de aceptaci\'on PA-02\\
\bottomrule
\end{tabularx}

\newpage
\subsection*{RNF-01 --- Tiempo de respuesta}
\begin{tabularx}{\linewidth}{@{}L{4.1cm} X@{}}
\toprule
\hd{Atributo} & \hd{Valor}\\
\midrule
Identificador & RNF-01\\
\rowcolor{rowgray} Tipo & No funcional\\
Nombre & Tiempo de registro de la cita\\
\rowcolor{rowgray} Caracter\'istica ISO/IEC 25010:2023 & Eficiencia de desempe\~no\\
Requisito & Desde que el paciente solicita agendar una cita hasta que el sistema presenta el resultado de la solicitud, el tiempo de respuesta deber\'a ser inferior a 3 segundos.\\
\rowcolor{rowgray} M\'etrica & Tiempo de respuesta\\
Umbral & $< 3$ segundos\\
\rowcolor{rowgray} Prioridad & Alta\\
Estado & En revisi\'on\\
\rowcolor{rowgray} Fuente & Caso pr\'actico: registrar la cita con rapidez\\
Estabilidad & Alta\\
\rowcolor{rowgray} Versi\'on & 1.0\\
M\'etodo de verificaci\'on & Prueba de rendimiento PA-03\\
\bottomrule
\end{tabularx}

\subsection*{RNF-02 --- Protecci\'on de datos personales}
\begin{tabularx}{\linewidth}{@{}L{4.1cm} X@{}}
\toprule
\hd{Atributo} & \hd{Valor}\\
\midrule
Identificador & RNF-02\\
\rowcolor{rowgray} Tipo & No funcional\\
Nombre & Protecci\'on de datos personales\\
\rowcolor{rowgray} Caracter\'istica ISO/IEC 25010:2023 & Seguridad\\
Requisito & El sistema deber\'a impedir que usuarios no autorizados visualicen la c\'edula, nombre y correo de los pacientes.\\
\rowcolor{rowgray} M\'etrica & Porcentaje de accesos no autorizados que permiten visualizar datos personales\\
Umbral & 0\%\\
\rowcolor{rowgray} Prioridad & Alta\\
Estado & En revisi\'on\\
\rowcolor{rowgray} Fuente & Caso pr\'actico: protecci\'on de datos personales del paciente\\
Estabilidad & Alta\\
\rowcolor{rowgray} Versi\'on & 1.0\\
M\'etodo de verificaci\'on & Prueba de seguridad PA-04\\
\bottomrule
\end{tabularx}

\newpage


\begin{tarea}{P6. Priorizaci\'on MoSCoW \hfill (5 pts)}
Clasifique los 4 requisitos de P5 en las categor\'ias \textbf{MoSCoW} (\textit{Must / Should / Could / Won't have this time}) y \textbf{justifique} brevemente cada asignaci\'on en funci\'on del valor y la viabilidad. Al menos un requisito debe ser \textit{Must have}.
\end{tarea}

\begin{tabularx}{\linewidth}{@{}L{3.5cm} C{2.4cm} X@{}}
\toprule
\hd{Requisito} & \hd{Prioridad MoSCoW} & \hd{Justificaci\'on}\\
\midrule
RF-01 --- Registrar cita & Must Have & Es indispensable para el funcionamiento del sistema, ya que permite registrar las citas solicitadas por los pacientes.\\
\rowcolor{rowgray}
RF-02 --- Verificar disponibilidad del cupo & Must Have & Es necesario para evitar la confirmaci\'on de citas en horarios que no tengan disponibilidad.\\
RNF-01 --- Tiempo de respuesta & Should Have & Es importante que el registro de una cita se realice en menos de 3 segundos para ofrecer un servicio r\'apido, aunque el sistema podr\'ia continuar funcionando si temporalmente supera ese tiempo.\\
\rowcolor{rowgray}
RNF-02 --- Protecci\'on de datos personales & Must Have & Es indispensable proteger la c\'edula, nombre y correo de los pacientes e impedir el acceso de usuarios no autorizados.\\
\bottomrule
\end{tabularx}

\newpage


\begin{tarea}{P7. Validaci\'on por inspecci\'on (lista de comprobaci\'on 29148) \hfill (8 pts)}
Aplique a los requisitos de P5 una \textbf{lista de comprobaci\'on} derivada de las propiedades de \textit{ISO/IEC/IEEE} 29148 (no ambiguo, completo, singular, verificable, factible). \textbf{Registre los defectos detectados sin corregirlos en el momento} (separar identificaci\'on de correcci\'on) y, en un segundo paso, presente el \textbf{retrabajo} con los requisitos corregidos.
\end{tarea}

\subsection*{1. Inspecci\'on de requisitos}
\scriptsize
\begin{tabularx}{\linewidth}{@{}L{2.7cm} C{1.45cm} C{1.3cm} C{1.3cm} C{1.45cm} C{1.25cm} X@{}}
\toprule
\hd{Requisito} & \hd{No ambiguo} & \hd{Completo} & \hd{Singular} & \hd{Verificable} & \hd{Factible} & \hd{Defecto}\\
\midrule
RF-01 Registrar cita & S\'i & S\'i & S\'i & S\'i & S\'i & Ninguno\\
\rowcolor{rowgray}
RF-02 Verificar disponibilidad del cupo & S\'i & S\'i & S\'i & S\'i & S\'i & Ninguno\\
RNF-01 Tiempo de respuesta & No & S\'i & S\'i & S\'i & S\'i & ``Resultado de la solicitud'' no especifica claramente cu\'al resultado.\\
\rowcolor{rowgray}
RNF-02 Protecci\'on de datos personales & No & S\'i & S\'i & S\'i & S\'i & ``Usuario no autorizado'' no define con precisi\'on qui\'en carece de autorizaci\'on.\\
\bottomrule
\end{tabularx}
\normalsize

\subsection*{2. Registro de defectos --- antes de corregir}
\textbf{DEF-01}\\
\textbf{Requisito:} RNF-01.\\
\textbf{Propiedad afectada:} No ambiguo.\\
La expresi\'on ``resultado de la solicitud'' puede interpretarse de diferentes maneras, ya que no especifica si el resultado corresponde a la confirmaci\'on de la cita o a la notificaci\'on de no disponibilidad.

\textbf{DEF-02}\\
\textbf{Requisito:} RNF-02.\\
\textbf{Propiedad afectada:} No ambiguo.\\
La expresi\'on ``usuarios no autorizados'' no especifica claramente qu\'e condici\'on determina que un usuario no posea autorizaci\'on para consultar los datos personales del paciente.

\subsection*{3. Retrabajo de requisitos}
\textbf{RNF-01 corregido --- Versi\'on 1.1}\\
Desde que el paciente solicita agendar una cita hasta que el sistema muestra la confirmaci\'on de la cita o la notificaci\'on de no disponibilidad del cupo, el tiempo de respuesta deber\'a ser inferior a 3 segundos.\\
\textbf{M\'etrica:} Tiempo de respuesta.\\
\textbf{Umbral:} $< 3$ segundos.

\textbf{RNF-02 corregido --- Versi\'on 1.1}\\
El sistema deber\'a impedir que los usuarios que no posean permisos de consulta de datos personales visualicen la c\'edula, nombre y correo de los pacientes.\\
\textbf{M\'etrica:} Porcentaje de accesos sin permiso que permiten visualizar datos personales.\\
\textbf{Umbral:} 0\%.

\subsection*{4. Reinspecci\'on despu\'es del retrabajo}
\small
\begin{tabularx}{\linewidth}{@{}L{3.1cm} C{1.9cm} C{1.7cm} C{1.7cm} C{1.8cm} C{1.5cm}@{}}
\toprule
\hd{Requisito} & \hd{No ambiguo} & \hd{Completo} & \hd{Singular} & \hd{Verificable} & \hd{Factible}\\
\midrule
RF-01 v1.0 & S\'i & S\'i & S\'i & S\'i & S\'i\\
\rowcolor{rowgray}
RF-02 v1.0 & S\'i & S\'i & S\'i & S\'i & S\'i\\
RNF-01 v1.1 & S\'i & S\'i & S\'i & S\'i & S\'i\\
\rowcolor{rowgray}
RNF-02 v1.1 & S\'i & S\'i & S\'i & S\'i & S\'i\\
\bottomrule
\end{tabularx}
\normalsize

\newpage


\begin{tarea}{P8. Pruebas de aceptaci\'on trazadas \hfill (6 pts)}
Defina \textbf{una prueba de aceptaci\'on por requisito} (4 en total), indicando su \textbf{tipo} (funcional, no funcional o regresi\'on), sus \textbf{entradas}, el \textbf{criterio de \'exito} y el \textbf{requisito al que traza}. Cada prueba debe enlazar expl\'icitamente con al menos un requisito de P5.
\end{tarea}

\begin{tarea}{P9. Matriz de trazabilidad \hfill (6 pts)}
Construya la \textbf{matriz de trazabilidad} que cruce los requisitos con: su \textbf{fuente} (traza hacia atr\'as) y sus \textbf{modelos y pruebas de aceptaci\'on} (traza hacia adelante). Marque tambi\'en al menos una dependencia \textbf{horizontal} entre requisitos e \textbf{identifique} cualquier requisito hu\'erfano (sin prueba o sin modelo asociado).
\end{tarea}

\begin{tarea}{P10. Gesti\'on del cambio y l\'inea base \hfill (4 pts)}
Simule \textbf{una solicitud de cambio} sobre un requisito: reg\'istrela, \textbf{analice su impacto} usando la matriz de P9, indique la \textbf{decisi\'on del CCB} (aprobada/rechazada/diferida) y \textbf{actualice la versi\'on} del requisito. Finalmente, declare una \textbf{l\'inea base}: liste la \textbf{configuraci\'on} (versiones coherentes de ERS, modelos, matriz y pruebas) que se congela.
\end{tarea}

% #####################################################################
\section{Rúbrica de evaluaci\'on}

La calificaci\'on total es de \textbf{100 puntos}, distribuidos en \textbf{30\%} para el cuestionario rendido en el LMS (SGA) y \textbf{70\%} para el desarrollo de la prueba pr\'actica. Los \textbf{criterios de piso} de la Secci\'on~1 se aplican \emph{antes} de sumar puntos: si alguno se incumple, la calificaci\'on final es \textbf{CERO}, con independencia de los puntos logrados en cada criterio.

\begin{concepto}{Escala de niveles de logro (se aplica a cada criterio de la prueba pr\'actica)}
\textbf{Logrado (100\% del peso):} el producto cumple \'integramente lo solicitado, es correcto y consistente. \quad \textbf{Parcial (60\%):} cumple lo esencial pero con omisiones o errores menores. \quad \textbf{Insuficiente (30\%):} entrega incompleta o con errores conceptuales relevantes. \quad \textbf{No entregado (0\%):} ausente o no evaluable.
\end{concepto}

% ── Componente SGA 30% ──
\begin{table}[htbp]
	\centering
	\caption{Componente 1 --- Cuestionario en el LMS (SGA). Peso: 30 puntos.}
	\small
	\begin{tabularx}{\linewidth}{@{}L{4.6cm} C{1.6cm} X@{}}
		\toprule
		\hd{Criterio} & \hd{Peso} & \hd{Evidencia requerida}\\
		\midrule
		Nota obtenida en el cuestionario del SGA & 20 pts & Puntaje autocalificado por el LMS, proporcional al porcentaje de aciertos.\\
		\rowcolor{rowgray}
		Captura del \emph{resumen} del cuestionario en la car\'atula & 5 pts & Imagen legible del resumen de resultados incrustada en la portada del PDF.\\
		Captura de la \emph{evaluaci\'on} (intento) en la car\'atula & 5 pts & Imagen legible del intento/evaluaci\'on incrustada en la portada del PDF.\\
		\bottomrule
	\end{tabularx}
\end{table}

% ── Componente practica 70% (landscape) ──
\begin{landscape}
\begin{table}[htbp]
	\centering
	\caption{Componente 2 --- Desarrollo de la prueba pr\'actica. Peso: 70 puntos. (Logrado 100\% / Parcial 60\% / Insuficiente 30\% / No entregado 0\% del peso indicado.)}
	\scriptsize
	\renewcommand{\arraystretch}{1.2}
	\begin{tabularx}{\linewidth}{@{}L{3.0cm} C{0.9cm} X X X@{}}
		\toprule
		\hd{Criterio (tarea)} & \hd{Peso} & \hd{Logrado (100\%)} & \hd{Parcial (60\%)} & \hd{Insuficiente (30\%)}\\
		\midrule
		P1. Diagrama de clases UML & 8 & Clases, atributos tipados, operaciones y multiplicidades correctas y coherentes con el caso. & Faltan operaciones o alguna multiplicidad; estructura b\'asica correcta. & Clases incompletas o sin multiplicidades; errores de modelado.\\
		\rowcolor{rowgray}
		P2. Diagrama de actividades UML & 8 & Decisi\'on de \textit{cupo}, ambos caminos, convergencia y nodo final bien modelados. & Flujo correcto pero sin decisi\'on o sin convergencia expl\'icita. & Secuencia lineal sin l\'ogica de control del caso.\\
		P3. M\'aquina de estados UML & 8 & Estados, eventos, guarda y estado final; superestado bien usado. & Estados y eventos correctos, sin guarda o sin estado final. & Estados incompletos o transiciones sin evento.\\
		\rowcolor{rowgray}
		P4. Consistencia entre perspectivas & 7 & Tabla cruzada completa; inconsistencia detectada, documentada y corregida. & Tabla parcial o inconsistencia detectada pero no corregida. & Tabla ausente o sin an\'alisis de correspondencia.\\
		P5. Requisitos con esquema de atributos & 10 & 4 requisitos (2 F, 2 NF) con caracter\'istica 25010, m\'etrica, umbral y atributos completos. & Requisitos correctos con atributos incompletos o sin umbral. & Requisitos ambiguos o sin esquema de atributos.\\
		\rowcolor{rowgray}
		P6. Priorizaci\'on MoSCoW & 5 & Cuatro requisitos clasificados y justificados; al menos un \textit{Must}. & Clasificaci\'on correcta con justificaci\'on d\'ebil. & Clasificaci\'on sin justificar o mal aplicada.\\
		P7. Inspecci\'on 29148 (defectos + retrabajo) & 8 & Checklist aplicada; defectos registrados sin corregir y retrabajo presentado. & Defectos detectados sin retrabajo, o correcci\'on mezclada con detecci\'on. & Sin checklist o sin evidencia de defectos.\\
		\rowcolor{rowgray}
		P8. Pruebas de aceptaci\'on trazadas & 6 & 4 pruebas con tipo, entradas, criterio de \'exito y traza al requisito. & Pruebas sin tipo o sin criterio de \'exito claro. & Pruebas sin traza a requisitos.\\
		P9. Matriz de trazabilidad & 6 & Trazas hacia atr\'as, adelante y horizontal; hu\'erfanos identificados. & Matriz parcial (solo un sentido de traza). & Matriz ausente o sin correspondencias reales.\\
		\rowcolor{rowgray}
		P10. Gesti\'on del cambio y l\'inea base & 4 & Solicitud, impacto, decisi\'on CCB, versi\'on y configuraci\'on/baseline declarada. & Cambio registrado sin impacto o sin baseline. & Sin proceso de cambio ni l\'inea base.\\
		\midrule
		\hg{Total del componente pr\'actico} & \hg{70} & \multicolumn{3}{c}{\cellcolor{lightgray}\small\color{darkblue} Suma de puntos logrados por nivel en cada criterio.}\\
		\bottomrule
	\end{tabularx}
\end{table}
\end{landscape}

\begin{table}[htbp]
	\centering
	\caption{Resumen de ponderaci\'on final.}
	\small
	\begin{tabularx}{\linewidth}{@{}X C{2.4cm} C{3.0cm}@{}}
		\toprule
		\hd{Componente} & \hd{Peso} & \hd{Condici\'on}\\
		\midrule
		Cuestionario en el LMS (SGA) + capturas en car\'atula & 30 pts & Sujeto a criterios de piso\\
		\rowcolor{rowgray}
		Desarrollo de la prueba pr\'actica (P1--P10) & 70 pts & Sujeto a criterios de piso\\
		\midrule
		\hg{Total} & \hg{100 pts} & \hg{CERO si se incumple un criterio de piso}\\
		\bottomrule
	\end{tabularx}
\end{table}

% #####################################################################
\section{Marco normativo de referencia}

Esta prueba operacionaliza las propiedades de calidad de la norma \textit{ISO/IEC/IEEE} 29148:2018 \citep{iso29148_2018} y el modelo de calidad de producto \textit{ISO/IEC} 25010:2023 \citep{iso25010_2023}, junto con las t\'ecnicas de validaci\'on y gesti\'on descritas por Pohl \citep{pohl_2010}, la inspecci\'on formal de Fagan \citep{fagan_1976} y las buenas pr\'acticas recogidas en la gu\'ia SWEBOK v4 \citep{swebok_v4_2024}. Las notaciones de modelado siguen la especificaci\'on \textit{OMG UML} 2.5.1 \citep{omg_uml_2017}.

\vspace{4pt}
\renewcommand{\refname}{Referencias}
\bibliographystyle{plainnat}
\bibliography{referencias}

\end{document}
